% GENERATED FILE. Edit canonical lessons, then run npm run generate:books.
% canonical-source-hash: fnv1a64:044eccd5793ce01d
% canonical-lessons: GU-R15-u-matra-r4, GU-R15-chha-r4, GU-R15-ka-r4, GU-R15-nna-r4, GU-R15-sha-r4, GU-R15-name-exchange-r3

\chapter{Fourth Return: Script and Names}
\label{ch:gu-r4-script-names}
\hlchaptermodality{19}

\begin{chapteropening}
\textbf{By the end of this chapter:} \emph{I can retrieve five Gujarati doorway forms at their fourth spacing window, distinguish each from a nearby familiar form, and independently read and write a complete name exchange after a long interval.}

The chapter closes without adding an atom: the learner independently reads three shuffled introduction lines, then writes a complete name exchange from English cues after first closing exact R4 windows for the short-u sign, chha, ka, retroflex nna, and sha.
\end{chapteropening}

\section[short-u sign — long-ii sign]{The short-u hook returns at R4}

\label{lesson:GU-R15-u-matra-r4}



\begin{quote}
Say \emph{short u} once. Picture its small hook below a consonant before you look at any form.
\end{quote}

\subsection*{Your turn}
Read \textbf{\gu{ુ} \gu{ી} \gu{ુ}} once and score recognition \textbf{3/3}. The short-u sign introduced at position 30 has now travelled 61 lessons to position 91: this closes its R4 window. Repair only a missed sign.

\subsection*{Writing — from sound}
Keep the visible row covered. Write the three sound cues, then score writing \textbf{3/3} separately. A perfect reading score cannot replace the written one.

\begin{tcolorbox}[breakable,colback=teal!4,colframe=teal!35!black,title={Before you move on}]
Where does \textbf{\gu{ુ}} sit? (\textbf{Below its consonant.})
\end{tcolorbox}
\section[chha — sha]{Chha returns beside sha}

\label{lesson:GU-R15-chha-r4}



\begin{quote}
Picture \emph{chha}: two high lobes around a broad lower body.
\end{quote}

\subsection*{Your turn}
Read \textbf{\gu{શ} \gu{છ} \gu{છ}} once and score recognition \textbf{3/3}. Chha moved from position 31 to position 92, exactly 61 lessons, so its R4 window closes here.

\subsection*{Writing — from sound}
With every model covered, write the three cues and score writing \textbf{3/3} on its own. Repair only a missed form.

\begin{tcolorbox}[breakable,colback=teal!4,colframe=teal!35!black,title={Before you move on}]
Which form has two high lobes? (\textbf{\gu{છ}, chha.})
\end{tcolorbox}
\section[ka — ja]{Ka returns beside ja}

\label{lesson:GU-R15-ka-r4}



\begin{quote}
Picture \emph{ka}: its compact upper turn leads into the lower body.
\end{quote}

\subsection*{Your turn}
Read \textbf{\gu{ક} \gu{જ} \gu{ક}} once and score recognition \textbf{3/3}. Ka moved from position 32 to position 93, exactly 61 lessons, closing its R4 window.

\subsection*{Writing — from sound}
Cover every model, write the cues, and score writing \textbf{3/3} independently. Repair only a missed form.

\begin{tcolorbox}[breakable,colback=teal!4,colframe=teal!35!black,title={Before you move on}]
Write \emph{ka} once more without looking. (\textbf{\gu{ક}})
\end{tcolorbox}
\section[retroflex nna — na]{Retroflex nna returns beside na}

\label{lesson:GU-R15-nna-r4}



\begin{quote}
Say \emph{retroflex nna} and picture \textbf{\gu{ણ}} before looking.
\end{quote}

\subsection*{Your turn}
Read \textbf{\gu{ન} \gu{ણ} \gu{ણ}} once and score recognition \textbf{3/3}. Retroflex nna moved from position 33 to position 94, exactly 61 lessons, closing its R4 window.

\subsection*{Writing — from sound}
Cover every model, write the three cues, and score writing \textbf{3/3} separately. Repair only a missed form.

\begin{tcolorbox}[breakable,colback=teal!4,colframe=teal!35!black,title={Before you move on}]
Which form answers the retroflex cue? (\textbf{\gu{ણ}})
\end{tcolorbox}
\section[sha — sa]{Sha returns beside sa}

\label{lesson:GU-R15-sha-r4}



\begin{quote}
Say \emph{sha}, then picture \textbf{\gu{શ}} without drawing it yet.
\end{quote}

\subsection*{Your turn}
Read \textbf{\gu{શ} \gu{સ} \gu{શ}} once and score recognition \textbf{3/3}. Sha moved from position 34 to position 95, exactly 61 lessons, so its R4 window closes here.

\subsection*{Writing — from sound}
Cover every model, write the cues, and score writing \textbf{3/3} independently. Repair only a missed form.

\begin{tcolorbox}[breakable,colback=teal!4,colframe=teal!35!black,title={Before you move on}]
Write \emph{sha} once more from sound. (\textbf{\gu{શ}})
\end{tcolorbox}
\section[Practice]{The name exchange returns after a long interval}

\label{lesson:GU-R15-name-exchange-r3}



\begin{quote}
With the page covered, say the three pieces for \emph{my name is}: \emph{mārũ — nām — chhe}.
\end{quote}

\subsection*{Your turn}
Read the three lines once and score meaning retrieval \textbf{3/3}. Position 35 was already a zero-atom checkpoint, so position 96 uses the open slot to close R3 name-exchange debt rather than pretending to retrieve a nonexistent new atom.

\subsection*{Writing — from sound}
Cover every Gujarati model and write the three English-cued lines. Score written production \textbf{3/3} separately from reading; repair only a missed line.

\begin{tcolorbox}[breakable,colback=teal!4,colframe=teal!35!black,title={Before you move on}]
Which \emph{you} belongs in this polite exchange? (\textbf{tame}.)
\end{tcolorbox}
